% Compile with XeLaTeX; the final report is one A4 page.
\documentclass[UTF8,fontset=none,zihao=5,a4paper]{ctexart}
\usepackage[left=22mm,right=22mm,top=19mm,bottom=19mm]{geometry}
\usepackage{fontspec,xcolor}
\setmainfont{DejaVu Serif}
\setsansfont{DejaVu Sans}
\setmonofont{DejaVu Sans Mono}[Scale=0.86]
\setCJKmainfont{AR PL SungtiL GB}[AutoFakeBold=2]
\setCJKsansfont{Droid Sans Fallback}[AutoFakeBold=2]
\setCJKmonofont{AR PL SungtiL GB}
\definecolor{ink}{HTML}{19364A}
\ctexset{section={format=\normalsize\sffamily\bfseries\color{ink},beforeskip=9pt,afterskip=4pt}}
\linespread{1.08}
\setlength{\parindent}{2em}
\setlength{\parskip}{3pt}
\setlength{\emergencystretch}{2em}
\pagestyle{empty}
\newcommand{\code}[1]{\texttt{\detokenize{#1}}}
\begin{document}
\begin{center}
  {\Large\sffamily\bfseries\color{ink}Lab 0: Booting 实验报告}\par
  \vspace{4pt}
  潘孝圆\quad 24300240128\quad 2026 年 10 月 16 日
\end{center}

本次修改文件为 \code{src/main.c}。实际执行顺序是：CPU 0 清零 BSS，唤醒其余核心并初始化打印服务，输出自己的问候，再放行其余核心打印。

\section*{任务 1：CPU 0 输出问候}
\textbf{思路与实现：}在 CPU 0 分支中，等 \code{uart_init()} 和 \code{printk_init()} 完成后，使用 \code{printk()} 输出 \code{Hello, world! (Core 0)} 并换行。打印位于放行其他核心之前，因此 Core 0 首先出现。

\textbf{交流摘要：}我问“ok写完了，接下来要干啥”。模型先解释 \code{printk()} 类似 C 中的 \code{printf()}，再带我找到对应 TODO，说明为什么要等打印服务初始化、为什么字符串末尾需要换行。我根据讲解手动补上这一处输出。

\section*{任务 2：其余核心输出各自编号}
\textbf{思路与实现：}其他核心先等待 \code{boot_secondary_cpus} 变为 \code{true}，经过原有 \code{arch_fence()} 后，再使用 \code{cpuid()} 取得本核编号，将其转换为 \code{int} 并配合 \texttt{\%d} 输出。三个核心并行执行，获取打印锁的顺序不固定，所以后三行可以交换顺序。

\textbf{交流摘要：}在同一轮讲解中，模型接着说明“当前核心编号”如何填入格式字符串，区分固定的 Core 0 输出与随核心变化的输出。我参照接口示例填写另一个 TODO，再运行并反馈终端结果，检查各核心是否都出现一次。

\section*{任务 3：初始化前清零 BSS}
\textbf{思路与实现：}BSS 中本实验使用的未显式初始化全局、静态对象需要零初值，启动时应保证这块内存已清零。复用链接脚本已有的 \code{edata}、\code{end}，用 \code{extern char edata[], end[];} 引用地址，再在 \code{smp_init()} 之前调用：
\begin{center}
\code{memset(edata, 0, (usize)end - (usize)edata);}
\end{center}
清零范围为 \code{[edata, end)}，第三个参数是两地址相差的字节数；本次范围覆盖 BSS，\code{linker.ld} 保持原样。

\textbf{交流摘要：}我先问“怎么找起始地址和结束边界啊”，又追问“为啥这么做啊”。模型用地址标记和逐字节填零的例子，解释 \code{extern}、\code{memset} 三个参数及边界不包含在范围内的含义。讨论后采用已有符号，我手动完成声明和清零调用。

\section*{运行结果与工具使用}
执行 \code{cmake --build build --target qemu} 后，最终验证输出顺序为 Core 0、Core 2、Core 1、Core 3，四核各出现一次。原先字符串中的小写 \code{core} 经检查改为 \code{Core}。

本次使用 Codex 辅助学习，模型可见标识为 GPT-6。以上对话为教学过程摘要：我不断追问原理和操作位置，模型逐步讲解并提供关键调用示例，我据此手动修改并运行验证；报告也由该工具辅助整理与 LaTeX 排版。
\end{document}
